% Generated draft: requires mathematical and source review.
\documentclass[11pt]{article}
\usepackage[margin=25mm]{geometry}
\usepackage{fontspec}
\setmainfont{Latin Modern Roman}
\usepackage{amsmath,amssymb,mathtools,unicode-math}
\setmathfont{Latin Modern Math}
\usepackage{enumitem,needspace,xurl,hyperref}
\setlength{\parindent}{0pt}
\setlength{\parskip}{5pt}
\setlength{\emergencystretch}{3em}
\newcommand{\N}{\mathbb{N}}
\newcommand{\Z}{\mathbb{Z}}
\newcommand{\Q}{\mathbb{Q}}
\newcommand{\R}{\mathbb{R}}
\newcommand{\C}{\mathbb{C}}
\newcommand{\F}{\mathbb{F}}
\newcommand{\A}{\mathbb{A}}
\newcommand{\PP}{\mathbb{P}}
\newcommand{\E}{\mathbb{E}}
\newcommand{\eps}{\varepsilon}
\newcommand{\dd}{\,\mathrm{d}}
\begin{document}
\section*{UnsolvedMath ID: 1100601}
\textbf{Source problem code:} catalogue entry
\subsection{Definitions and mathematical statement}
Let $p$ be a prime and $\mathbb F_p$ the field with $p$ elements. Under the finite-resolution convention, a group $H$ is of type $\mathrm{FP}$ over $\mathbb F_p$ if the trivial left $\mathbb F_p[H]$-module $\mathbb F_p$ has an exact projective resolution
\[
0\longrightarrow P_m\longrightarrow\cdots\longrightarrow P_0\longrightarrow\mathbb F_p\longrightarrow0
\]
for some finite $m$, with every $P_i$ finitely generated. A group is virtually of this type if it has a finite-index subgroup of this type. For $g\in G$, put
\[
C_G(g)=\{x\in G:xg=gx\}.
\]
\textbf{Question (Ian Leary).} For every prime $p$, every group $G$ virtually of type $\mathrm{FP}$ over $\mathbb F_p$, and every $g\in G$ with order exactly $p$, is $C_G(g)$ virtually of type $\mathrm{FP}$ over $\mathbb F_p$? Equivalently, must there exist $L\leq C_G(g)$ with $[C_G(g):L]<\infty$ having the stated finiteness property?

The record reports an affirmative answer from Hamilton's Theorem E, with the stronger centralizer conclusion for an arbitrary $p$-subgroup $P\leq G$, where $C_G(P)=\{x\in G:xy=yx\text{ for every }y\in P\}$. The supplied statement does not define $\mathrm{FP}$: if it instead intends $\mathrm{FP}_{\infty}$ (a resolution with finitely generated projective modules in every degree, with no finite-length requirement), that is a different convention and requires checking against the cited theorem.
\subsection{Short English statement}
If a group is virtually homologically finite over $\mathbb F_p$, must the centralizer of an element of order $p$ have the same virtual finiteness property? The supplied literature summary reports yes, and a stronger result for centralizers of $p$-subgroups.
\subsection{Sources}
\begin{itemize}
\item M. Bestvina, \emph{Questions in Geometric Group Theory} (2004), Question 6.1, PDF page 14, attributed to Ian Leary: \url{https://www.math.utah.edu/~bestvina/eprints/questions-updated.pdf}; manual transcription, accessed 2026-07-29.
\item Hamilton, \emph{Journal of Algebra} 330 (2011), 1--21, Theorem E; Leary--Nucinkis, \emph{Inventiones Mathematicae} 151 (2003), Question 1. These are the supplied research-summary citations; article titles are not provided.
\item Local registry \texttt{data/unsolvedmath/problems.json}, permanent ID $1100601$, source code catalogue entry, \emph{Questions in Geometric Group Theory --- Q 6.1}. The record labels this solved and SOLVED-IN-LITERATURE; its source-list background reports partial progress and retains current-status and rights-review qualifications.
\end{itemize}
\end{document}
